% !TeX encoding = UTF-8
% !TeX TS-program = xelatex
% 通用 Beamer 模板主示例

\documentclass[hyperref, UTF8, CJK, aspectratio=169]{beamer}

% --------通用 Beamer 主题--------
\usetheme[
  ColorDisplay=Blue,
  Miniframes=follow,
  ContentMuticols=true,
  BIBMode=none,
  Background=false,
]{academic}

\usepackage{amsmath,amssymb}
\usepackage{booktabs}
\usepackage{graphicx}
\usepackage{tabularx}
\usepackage{pgfplots}
\usetikzlibrary{arrows.meta}
\pgfplotsset{compat=1.18,every axis/.append style={font=\academicfont{chart-label}}}
\pgfplotscreateplotcyclelist{academicChartCycle}{%
  {ChartPrimaryC,solid,thick},%
  {ChartSecondaryC,dashed,thick},%
  {ChartContrastC,dashdotted,thick},%
  {ChartFourthC,densely dashdotted,thick}%
}
\pgfplotsset{%
  academic chart/.style={%
    cycle list name=academicChartCycle,
    axis background/.style={fill=white},
    axis line style={draw=ChartAxisC},
    tick label style={color=ChartAxisC},
    label style={color=ChartAxisC},
    grid style={draw=ChartGridC},
    legend style={draw=BorderC,fill=white,text=NormalTextC}
  }%
}

% 调色板示例页使用的色块标签。
\newcommand{\academicswatch}[3]{%
  \noindent\tikz[baseline=-.5ex]\filldraw[draw=BorderC,fill=#1]
    (0,0) rectangle (.85em,.65em);%
  \hspace{.35em}{\academicfont{label}\texttt{#1}}\enspace{\academicfont{label} #3}\enspace{\academicfont{caption} #2}\par%
}

% --------封面信息--------
\title[汇报题目]{汇报题目}
\subtitle{ 副标题（可选）}
\author[ 作者姓名]{作者姓名}
\institute{单位名称}
\date{\today}

\begin{document}

% --------主题与导航--------
\section{主题与导航}
\subsection{章节目录与页面布局}

\begin{frame}{主题结构与目录}
  \begin{columns}[T,onlytextwidth]
    \begin{column}{.48\textwidth}
      \begin{ResearchCard}{自动目录}
        主目录列出大章；顶部导航同步显示当前进度。
      \end{ResearchCard}
    \end{column}
    \begin{column}{.48\textwidth}
      \begin{ResearchCard}{页面导航}
        章节导航、页码和大章目录由主题统一管理。
      \end{ResearchCard}
    \end{column}
  \end{columns}
  \vspace{.5em}
  \begin{itemize}
    \item 在 \texttt{main.tex} 中填写标题，章节名称自动生成目录。
    \item 封面装饰使用矢量资源 \texttt{resources/cover-background.pdf}。
  \end{itemize}
\end{frame}

\subsection{基础区块与列表}
\begin{frame}{基础 Beamer 组件}
  \begin{columns}[T,onlytextwidth]
    \begin{column}{.32\textwidth}
      \begin{ResearchCard}{普通区块}
        一般说明。
      \end{ResearchCard}
    \end{column}
    \begin{column}{.32\textwidth}
      \begin{ResearchCard}{示例区块}
        方法案例。
      \end{ResearchCard}
    \end{column}
    \begin{column}{.32\textwidth}
      \begin{ResearchCard}{强调区块}
        重要提醒。
      \end{ResearchCard}
    \end{column}
  \end{columns}
  \begin{enumerate}
    \item 资源放入 \texttt{resources/} 等目录。
    \item 使用 XeLaTeX；文献模块再启用 Biber。
  \end{enumerate}
\end{frame}

\subsection{统一字号与强调}
\begin{frame}{统一字号接口与局部强调}
  \begin{columns}[T,onlytextwidth]
    \begin{column}{.48\textwidth}
      {\academicstandardtext
      \begin{ResearchCard}{统一正文接口}
        深灰蓝正文；\alert{砖橙一级重点}，\academicemphblue{工程蓝二级重点}。
      \end{ResearchCard}}
    \end{column}
    \begin{column}{.48\textwidth}
      {\academicfocustext
      \begin{ResearchCard}{统一标题字号}
        页面标题和卡片标题均通过统一接口设置。
      \end{ResearchCard}}
    \end{column}
  \end{columns}
  \vspace{.35em}
  \centering 正文 23pt，页面标题与卡片标题 25pt；所有组件字号集中配置。
\end{frame}

\begin{frame}{ResearchCard：完整卡片与可选编号}
  \begin{ResearchCard}[1]{完整卡片：有数字编号}
    卡片样式统一；整数编号是可选的显示属性。
  \end{ResearchCard}
  \vspace{.35em}
  \begin{ResearchCard}{完整卡片：无数字编号}
    不提供整数序号时，不绘制圆形编号。
  \end{ResearchCard}
\end{frame}

% --------统一卡片：按显示布局展示--------
\section{统一研究卡片}
\subsection{双栏卡片}
\begin{frame}{双栏卡片：文字内容}
  \begin{columns}[T,onlytextwidth]
    \begin{column}{.48\textwidth}
      \begin{ResearchCard}{定义：示例概念}
        此处填写核心术语的定义。
      \end{ResearchCard}
    \end{column}
    \begin{column}{.48\textwidth}
      \begin{ResearchCard}{条件：适用条件}
        此处填写方法或结论成立所需的前提条件。
      \end{ResearchCard}
    \end{column}
  \end{columns}
\end{frame}

\subsection{双栏卡片：不同标题内容}
\begin{frame}{双栏卡片：不同标题内容}
  \begin{columns}[T,onlytextwidth]
    \begin{column}{.48\textwidth}
      \begin{ResearchCard}{定理：定理示例}
        此处填写定理陈述及其适用前提。
      \end{ResearchCard}
    \end{column}
    \begin{column}{.48\textwidth}
      \begin{ResearchCard}{符号约定}
        明确变量含义、方向约定和计量单位。
      \end{ResearchCard}
    \end{column}
  \end{columns}
\end{frame}

\subsection{双栏与通栏组合}
\begin{frame}{双栏与通栏卡片}
  \begin{columns}[T,onlytextwidth]
    \begin{column}{.48\textwidth}
      \begin{ResearchCard}{公式：通用函数关系}
        \[
          y = f(x;\theta).
        \]
      \end{ResearchCard}
    \end{column}
    \begin{column}{.48\textwidth}
      \begin{ResearchCard}{示例：示例案例}
        此处填写代表性算例。
      \end{ResearchCard}
    \end{column}
  \end{columns}
  \begin{ResearchCard}{结论：结论示例}
    概括主要发现与适用范围。
  \end{ResearchCard}
\end{frame}

\subsection{单栏卡片}
\begin{frame}{单栏卡片：研究问题}
  \begin{ResearchCard}{问题：核心问题}
    此处填写本研究需要回答的核心问题。
  \end{ResearchCard}
  \vspace{.5em}
  \begin{itemize}
    \item 可在此处列出研究假设或待验证的问题。
    \item 问题表述应与研究目标保持一致。
  \end{itemize}
\end{frame}

% --------表格与程序模块--------
\section{表格与程序}
\subsection{表头与交替行}
\begin{frame}{表格配色模块}
  \begin{table}
    \centering
    \academicfont{table}
    \academictablebegin
    \begin{tabular}{lcc}
      \academictabletoprule
      \rowcolor{TableHeaderC}
      \academictableheader{项目} & \academictableheader{设置} & \academictableheader{结果} \\
      \academictableheadrule
      条目 A & 待填写 & 待填写 \\
      \academictablebodyrule
      条目 B & 待填写 & 待填写 \\
      \academictablebodyrule
      条目 C & 待填写 & 待填写 \\
      \academictablebottomrule
    \end{tabular}
    \academictableend
  \end{table}
  \vspace{.5em}
  \begin{ResearchCard}{表格颜色}
    表头、正文与规则线均由主题颜色变量统一控制。
  \end{ResearchCard}
\end{frame}

\subsection{代码与重点行高亮}
\begin{frame}[fragile]{代码高亮模块}
  \begin{CodePanel}{数据处理示例}{Python}[2,3]
values = [1.0, 2.0, 3.0]
mean = sum(values) / len(values)
print(f"mean: {mean:.2f}")
  \end{CodePanel}
  \vspace{.35em}
  \begin{CodeLine}{峰值响应提取}{Python}
peak = np.max(np.abs(response))
  \end{CodeLine}
\end{frame}

% --------图表与色板--------
\section{科研图表与色板}
\subsection{图表主线与对比线}
\begin{frame}{图表配色模块}
  \centering
  \begin{tikzpicture}[x=\academicdim{1cm},y=\academicdim{1cm}]
    \begin{axis}[
      academic chart,
      width=.86\linewidth,
      height=.54\textheight,
      xlabel={自变量 $x$},
      ylabel={输出量 $y$},
      legend columns=3,
      legend to name=academic-main-chart-legend,
      clip mode=individual,
      grid=major,
    ]
      \addplot coordinates {(0,0.10) (1,0.35) (2,0.55) (3,0.62) (4,0.48)};
      \node[anchor=south west,text=PrimaryC,font=\academicfont{caption}] at (axis cs:3,0.62) {示例点};
      \addlegendentry{方案 A}
      \addplot coordinates {(0,0.08) (1,0.25) (2,0.42) (3,0.50) (4,0.40)};
      \addlegendentry{方案 B}
      \addplot coordinates {(0,0.03) (1,0.18) (2,0.30) (3,0.38) (4,0.32)};
      \addlegendentry{方案 C}
    \end{axis}
  \end{tikzpicture}
  \par\smallskip
  \pgfplotslegendfromname{academic-main-chart-legend}
  \par\smallskip
  {\academicfont{caption}\color{AuxiliaryC}示意数据仅用于展示图表样式，使用时请替换为实际数据。}
\end{frame}

\subsection{研究流程图}
\begin{frame}{研究流程图模块}
  \centering
  \begin{tikzpicture}[x=\academicdim{1cm},y=\academicdim{1cm},
    flow step/.style={
      draw=AuxiliaryC,
      rounded corners=\academicdim{1.4mm},
      line width=\academicdim{.75pt},
      minimum width=\academicdim{2.10cm},
      minimum height=\academicdim{1.42cm},
      text width=\academicdim{1.88cm},
      align=center,
      inner sep=\academicdim{3.5pt},
      font=\academicfont{diagram},
      text=NormalTextC
    },
    flow arrow/.style={-{Stealth[length=\academicdim{2mm}]},draw=IntersperseC,line width=\academicdim{1pt}},
    flow feedback/.style={-{Stealth[length=\academicdim{1.8mm}]},draw=AuxiliaryC,line width=\academicdim{.8pt}}
  ]
    \node[flow step,fill=BlockQuestionC] (question) at (-4.8,0) {\textcolor{PrimaryC}{\bfseries 01\quad 研究目标}\\[-\academicdim{1pt}]\academicfont{caption} 明确目标与研究范围};
    \node[flow step,fill=BlockDefinitionC] (theory) at (-2.4,0) {\textcolor{PrimaryC}{\bfseries 02\quad 研究方案}\\[-\academicdim{1pt}]\academicfont{caption} 确定假设与方法};
    \node[flow step,fill=BlockFormulaC] (model) at (0,0) {\textcolor{PrimaryC}{\bfseries 03\quad 分析实施}\\[-\academicdim{1pt}]\academicfont{caption} 开展分析或计算};
    \node[flow step,fill=BlockExampleC] (validation) at (2.4,0) {\textcolor{PrimaryC}{\bfseries 04\quad 结果验证}\\[-\academicdim{1pt}]\academicfont{caption} 检查结果与误差};
    \node[flow step,fill=BlockConclusionC] (conclusion) at (4.8,0) {\textcolor{PrimaryC}{\bfseries 05\quad 主要结论}\\[-\academicdim{1pt}]\academicfont{caption} 整理主要发现};

    \draw[flow arrow] (question.east) -- (theory.west);
    \draw[flow arrow] (theory.east) -- (model.west);
    \draw[flow arrow] (model.east) -- (validation.west);
    \draw[flow arrow] (validation.east) -- (conclusion.west);
    \draw[flow feedback] (validation.south) -- ++(0,-.48) -|
      node[pos=.56,below=\academicdim{1pt},fill=BackgroundC,inner sep=\academicdim{1.5pt},font=\academicfont{caption},text=NormalTextC]
      {检查后调整方法并复算} (model.south);
  \end{tikzpicture}

  \vspace{.7em}
  {\academicfont{caption}\color{AuxiliaryC}流程节点可按实际工作内容增删，并调整顺序与反馈关系。}
\end{frame}

\subsection{主题色与模块色总览}
\begin{frame}{完整配色变量总览}
  \begin{columns}[T,onlytextwidth]
    \begin{column}{.32\textwidth}
      \textbf{基础颜色}\par\smallskip
      \academicswatch{PrimaryC}{学术深蓝}{\#173B75}
      \academicswatch{AuxiliaryC}{工程蓝}{\#235F9E}
      \academicswatch{SecondaryAuxiliaryC}{海松绿}{\#146F68}
      \academicswatch{IntersperseC}{砖橙色}{\#B45929}
      \academicswatch{NormalTextC}{正文色}{\#243B55}
      \academicswatch{AlertedTextC}{重要数据}{\#B45929}
      \academicswatch{MutedTextC}{次要文字}{\#52606C}
      \academicswatch{BackgroundC}{背景色}{\#FFFFFF}
      \academicswatch{BorderC}{分隔线}{\#7895AF}
    \end{column}
    \begin{column}{.34\textwidth}
      \textbf{内容模块}\par\smallskip
      \academicswatch{BlockNormalBgC}{普通/定理}{\#FFFFFF}
      \academicswatch{BlockWarmBgC}{结论白底}{\#FFFFFF}
      \academicswatch{BlockCoolBgC}{理论/问题白底}{\#FFFFFF}
      \academicswatch{BlockDefinitionC}{定义}{\#FFFFFF}
      \academicswatch{BlockTheoremC}{定理}{\#FFFFFF}
      \academicswatch{BlockLemmaC}{引理}{\#FFFFFF}
      \academicswatch{BlockConditionC}{条件}{\#FFFFFF}
      \academicswatch{BlockExampleC}{示例}{\#FFFFFF}
      \academicswatch{BlockFormulaC}{公式}{\#FFFFFF}
      \academicswatch{BlockConclusionC}{结论}{\#FFFFFF}
      \academicswatch{BlockQuestionC}{问题}{\#FFFFFF}
    \end{column}
    \begin{column}{.34\textwidth}
      \textbf{表格、图表与代码}\par\smallskip
      \academicswatch{TableHeaderC}{表头背景}{\#F2F6FB}
      \academicswatch{TableRuleC}{顶部/底部线}{\#235F9E}
      \academicswatch{TableRowRuleC}{表头分界线}{\#7294B5}
      \academicswatch{ChartPrimaryC}{主曲线}{\#235F9E}
      \academicswatch{ChartSecondaryC}{第二曲线}{\#B45929}
      \academicswatch{ChartContrastC}{第三曲线}{\#146F68}
      \academicswatch{ChartFourthC}{第四曲线}{\#70528B}
      \academicswatch{ChartAxisC}{坐标轴}{\#243B55}
      \academicswatch{ChartGridC}{网格线}{\#A5B4C2}
      \academicswatch{CodeCommentC}{代码注释}{\#146F68}
      \academicswatch{HighlightCodeLineC}{重点行}{\#E6F0FA}
    \end{column}
  \end{columns}
\end{frame}

% --------通用研究方案示例--------
\section{研究方案示例}
\subsection{模型示意与参数设置}
\begin{frame}{模型示意与参数设置}
  \begin{columns}[T,onlytextwidth]
    \begin{column}{.52\textwidth}
      \centering
      \begin{tikzpicture}[x=\academicdim{1cm},y=\academicdim{1cm},
        model box/.style={draw=AuxiliaryC,line width=\academicdim{.8pt},
          rounded corners=\academicdim{1mm},minimum height=\academicdim{1cm},text width=\academicdim{1.35cm},
          align=center,font=\academicfont{caption},text=NormalTextC},
        model arrow/.style={-{Stealth[length=\academicdim{1.8mm}]},draw=IntersperseC,line width=\academicdim{.8pt}}
      ]
        \node[model box,fill=BlockConditionC] (input) at (-2.15,0) {输入参数};
        \node[model box,fill=BlockDefinitionC,minimum width=\academicdim{1.8cm}] (model) at (0,0) {研究对象\\模型};
        \node[model box,fill=BlockExampleC] (output) at (2.15,0) {输出结果};
        \draw[model arrow] (input.east) -- (model.west);
        \draw[model arrow] (model.east) -- (output.west);
      \end{tikzpicture}
      \par\vspace{.2em}
      {\academicfont{caption}\color{AuxiliaryC}通用模型示意，可替换为理论模型、计算模型或有限元网格。}
    \end{column}
    \begin{column}{.45\textwidth}
      \begin{ResearchCard}{假设：示例假设}
        此处填写模型假设、适用范围和需要说明的限制。
      \end{ResearchCard}
      \vspace{.35em}
      \academicfont{table}
      \academictablebegin
      \begin{tabular}{@{}ll@{}}
        \academictabletoprule
        \rowcolor{TableHeaderC}
        \academictableheader{参数项} & \academictableheader{符号 / 来源} \\
        \academictableheadrule
        参数 A & $p_A$ / 数据来源 \\
        \academictablebodyrule
        参数 B & $p_B$ / 数据来源 \\
        \academictablebodyrule
        控制变量 & $c$ / 工况设置 \\
        \academictablebottomrule
      \end{tabular}
      \academictableend
    \end{column}
  \end{columns}
\end{frame}

\subsection{试验验证与参数敏感性}
\begin{frame}{试验验证与参数敏感性}
  \begin{columns}[T,onlytextwidth]
    \begin{column}{.48\textwidth}
      \centering
      \begin{tikzpicture}[x=\academicdim{1cm},y=\academicdim{1cm},
        rig box/.style={draw=AuxiliaryC,rounded corners=\academicdim{1mm},
          fill=BlockConditionC,minimum height=\academicdim{8mm},text=NormalTextC,
          font=\academicfont{caption},align=center},
        rig arrow/.style={-{Stealth[length=\academicdim{1.8mm}]},draw=IntersperseC,line width=\academicdim{.8pt}}
      ]
        \node[rig box,minimum width=\academicdim{1.5cm}] (drive) at (-1.8,.45) {输入 / 加载};
        \node[draw=PrimaryC,line width=\academicdim{1pt},fill=BlockDefinitionC,circle,
          minimum size=\academicdim{1.12cm},font=\academicfont{caption},text=PrimaryC] (object) at (0,.45) {研究对象};
        \node[rig box,minimum width=\academicdim{1.5cm}] (acq) at (1.85,.45) {测量 / 采集};
        \draw[rig arrow] (drive.east) -- (object.west);
        \draw[rig arrow] (object.east) -- (acq.west);
        \draw[AuxiliaryC,line width=\academicdim{2pt}] (-.68,-.48) -- (.68,-.48);
        \node[font=\academicfont{caption},text=NormalTextC] at (0,-.7) {试验平台};
        \draw[rig arrow] (.32,-.36) -- (.32,.05)
          node[midway,right=\academicdim{1pt},text=NormalTextC,font=\academicfont{caption}] {$u$};
      \end{tikzpicture}
      \par\vspace{.2em}
      {\academicfont{caption}\color{AuxiliaryC}装置示意可替换为实际设备图、测点或实验流程。}
      \vspace{.4em}
      \academictablebegin
      \begin{tabular}{@{}lc@{}}
        \academictabletoprule
        \rowcolor{TableHeaderC}
        \academictableheader{评价指标} & \academictableheader{实测 / 计算} \\
        \academictableheadrule
        指标 A & 待填写 \\
        \academictablebodyrule
        指标 B & 待填写 \\
        \academictablebodyrule
        指标 C & 待填写 \\
        \academictablebottomrule
      \end{tabular}
      \academictableend
    \end{column}
    \begin{column}{.49\textwidth}
      \centering
      \begin{tikzpicture}[x=\academicdim{1cm},y=\academicdim{1cm}]
        \begin{axis}[
          academic chart,
          width=\linewidth,
          height=.49\textheight,
          xlabel={参数变化率（\%）},
          ylabel={归一化输出},
          xmin=-22,xmax=22,ymin=.65,ymax=1.4,
          xtick={-20,-10,0,10,20},
          ytick={.7,.9,1.1,1.3},
          grid=major,
          title={敏感性分析示意},
          title style={font=\academicfont{chart-title},text=PrimaryC}
        ]
          \addplot+[ChartPrimaryC,mark=*,line width=\academicdim{1.1pt}]
            coordinates {(-20,.76) (-10,.88) (0,1.00) (10,1.14) (20,1.30)};
        \end{axis}
      \end{tikzpicture}
      \par\vspace{-.2em}
      {\academicfont{caption}\color{AuxiliaryC}示意数据仅用于展示敏感性分析图样式。}
    \end{column}
  \end{columns}
\end{frame}

% --------结束页--------
\section{使用建议}
\begin{frame}{开始编写你的汇报}
  \begin{ResearchCard}{推荐流程}
    先在主文件中组织章节，再填充对应模块，最后检查目录、图表和编译结果。
  \end{ResearchCard}
  \begin{enumerate}
    \item 在 \texttt{main.tex} 中修改标题和正文；
    \item 用 \texttt{section}/\texttt{subsection} 组织内容与目录；
    \item 将图表放入 \texttt{figure/} 或 \texttt{resources/}；
    \item 使用 XeLaTeX 编译两遍，让自动目录完整更新。
  \end{enumerate}
\end{frame}

% --------科研汇报结束页--------
\ClosingPage{汇报完毕}{感谢各位专家批评指正}

\end{document}
